\subsection*{Loops in Croqtile}

In C++, loops are typically constructed using \texttt{for}, \texttt{while}, and \texttt{do-while} blocks. These loops include conditions for termination and usually involve variable manipulations related to loop control. Programmers have the freedom to define exit conditions, loop stepping, and other aspects, allowing for flexible loop structures.

Croqtile uses bounded ranges for these loops. The range is defined once and then visited by \co@foreach@:

\begin{lstlisting}[language=croqtile,caption={A foreach loop}]
foreach index {
  // loop body
}
\end{lstlisting}

Here, \texttt{index} is a \emph{bounded variable} (introduced in the previous section). If \texttt{index} has an upper bound of \texttt{6}, the above code is equivalent to:

\begin{lstlisting}[language={},caption={The equivalent C++ loop}]
for (int index = 0; index < 6; ++index) {
  // loop body
}
\end{lstlisting}

This form suits tile traversal: the loop visits the valid coordinates in a declared range, and the same bounds can be used to derive each tile's shape.

In Croqtile, \emph{bounded variables} defined within a \texttt{parallel-by} statement are immutable and cannot be used in \texttt{foreach} statements. Instead, programmers should use the \texttt{with-in} statement to define the \emph{bounded variables} suitable for loops.
